\cotefile{162-5}
\scanpage{30}
9.\quad \deleted{\unsure{Partant} d'une} \uline{Opération d'une $\Theta$.Alg. de Lie}

Soit $\mathscr T$ une $\Theta$-Algèbre de Lie sur $X_0$, $M$ un $\underline O_{X_0}$-Module. On dit que \deleted{$\mathscr T$} $M$ a une $\mathscr T$-connexion si on fait ``\unsure{opérer|operer}'' $\mathscr T$ sur $M$ de façon à satisfaire à a) et b) plus haut.

Ceci définit alors un tenseur courbure
\[
K\in C^2(\mathscr T,\mathrm{End}_{\underline O_{X_0}}(M))
\]
en posant
\[
K(X,Y)=\Theta_{[X,Y]}^M-\unsure{[|(}\Theta_X^M,\Theta_Y^H]
\]
Le tenseur courbure est donc nul ssi on a
\[
\Theta_{[X,Y]}^M=[\Theta_X^M,\Theta_Y^H],
\]
alors la connexion est \uline{dite intégrable}. Si $\mathscr T$ provient de $\Omega$ \inserted{$(\mathscr T=\underline{\mathrm{Hom}}_{\underline O_{X_0}}(\Omega,\underline O_{X_0}))$}, et $M$ \deleted{\unsure{provient}} muni d'une $\Omega$-connexion $\delta_M^0$, d'où un tenseur courbure
\diagram{\begin{tikzcd}[ampersand replacement=\&,column sep=large]
M\arrow[r]\arrow[rr,bend right=22,no head,"K'"'] \& \bigwedge^2\Omega\otimes M\arrow[r] \& C^2(\mathscr T,M)
\end{tikzcd}}
\unsure{si conventions|j'écris maintenant|j'aimerais}
\margin{Attention signes !}
\[
K(X,Y)(u)=K'(u)(X,X).
\]

\uline{10}.\quad Si \inserted{$(\bigwedge^*\Omega,\delta)$ avec défini\unsure{.}} $\Omega$ loc. libre de t.f. ($\Longleftrightarrow\mathscr T$ \inserted{$\Theta$-alg. de Lie sur $\mathscr T$} loc. libre de t.f.) alors pour un Module $M$ sur $\underline O_{X_0}$,

\margin{29}
